% !TEX program = pdflatex
% ==========================================================
% PLANTILLA BEAMER PROFESIONAL
% Universidad Santo Tomás
% Luis Tulio González Alcaino
% ==========================================================

\newif\ifhandout
\handoutfalse
%\handouttrue

\ifhandout
\documentclass[aspectratio=169,10pt,handout]{beamer}
\else
\documentclass[aspectratio=169,10pt]{beamer}
\fi

% ----------------------------------------------------------
% PAQUETES
% ----------------------------------------------------------

\usepackage[spanish,es-noshorthands]{babel}
\usepackage[T1]{fontenc}
\usepackage[utf8]{inputenc}

\usepackage{lmodern}
\usepackage{amsmath,amsfonts,amssymb}
\usepackage{ragged2e}
\usepackage{enumitem}

\usepackage{tikz}
\usetikzlibrary{positioning}

\usepackage{fontawesome5}

% ----------------------------------------------------------
% COLORES INSTITUCIONALES
% ----------------------------------------------------------

\definecolor{USTBlue}{RGB}{0,70,127}
\definecolor{USTBlueDark}{RGB}{0,45,85}
\definecolor{USTTeal}{RGB}{0,153,117}
\definecolor{USTGray}{RGB}{120,120,120}
\definecolor{USTSoft}{RGB}{248,249,250}
\definecolor{USTText}{RGB}{35,40,45}
\definecolor{WarmGray}{RGB}{120,120,120}

% ----------------------------------------------------------
% TEMA BASE
% ----------------------------------------------------------

\usetheme{Madrid}
\usecolortheme[named=USTBlue]{structure}

% eliminar iconos navegación
\setbeamertemplate{navigation symbols}{}

% ----------------------------------------------------------
% COLORES DE PIE
% ----------------------------------------------------------

\setbeamercolor{author in head/foot}{bg=USTSoft,fg=USTGray}
\setbeamercolor{date in head/foot}{bg=USTSoft,fg=USTGray}
\setbeamercolor{title in head/foot}{bg=USTSoft,fg=USTGray}

% ----------------------------------------------------------
% PIE DE PÁGINA
% ----------------------------------------------------------

\setbeamertemplate{footline}{%
	\leavevmode%
	\hbox{%
		\begin{beamercolorbox}[wd=.85\paperwidth,ht=2.5ex,dp=1ex,leftskip=0.4cm]{author in head/foot}%
			\scriptsize
			\textbf{Fundamentos de Matemática Básica}
			\quad|\quad
			Agronomía
			\quad|\quad
			Universidad Santo Tomás Talca
		\end{beamercolorbox}%
		\begin{beamercolorbox}[wd=.15\paperwidth,ht=2.5ex,dp=1ex,rightskip=0.3cm plus1fil]{date in head/foot}%
			\scriptsize \insertframenumber/\inserttotalframenumber
		\end{beamercolorbox}%
	}%
}

% ----------------------------------------------------------
% BLOQUES PEDAGÓGICOS
% ----------------------------------------------------------

\setbeamercolor{block title}{bg=USTBlue!15,fg=USTBlueDark}
\setbeamercolor{block body}{bg=USTSoft,fg=black}

\setbeamercolor{block title example}{bg=USTTeal!20,fg=USTBlueDark}
\setbeamercolor{block body example}{bg=USTSoft,fg=black}

\setbeamercolor{block title alerted}{bg=red!15,fg=black}
\setbeamercolor{block body alerted}{bg=red!3,fg=black}

% ----------------------------------------------------------
% COMANDOS Y ENTORNOS AUXILIARES
% ----------------------------------------------------------

\newcommand{\destacar}[1]{\textcolor{USTBlueDark}{\textbf{#1}}}
\newcommand{\alerta}[1]{\textcolor{red!70!black}{\textbf{#1}}}
\newcommand{\pp}{\pause}

\newenvironment{metaBox}
{\begin{block}{Meta de aprendizaje}}
	{\end{block}}

\newenvironment{ideaBox}
{\begin{exampleblock}{Idea clave}}
	{\end{exampleblock}}

% ----------------------------------------------------------
% PORTADA PROFESIONAL
% ----------------------------------------------------------

% ----------------------------------------------------------
% PORTADA PROFESIONAL PREMIUM
% ----------------------------------------------------------
\newcommand{\PortadaProfesionalUST}{%
	\begin{frame}[plain]
		\begin{tikzpicture}[remember picture,overlay]
			
			% Fondo general
			\fill[white] (current page.south west) rectangle (current page.north east);
			
			% Franja superior principal
			\fill[USTBlueDark] (current page.north west) rectangle ([yshift=-1.35cm]current page.north east);
			
			% Línea de acento
			\fill[USTTeal] ([yshift=-1.35cm]current page.north west) rectangle ([yshift=-1.52cm]current page.north east);
			
			% Panel lateral suave
			\fill[USTSoft] (current page.south west) rectangle ([xshift=2.15cm]current page.north west);
			
			% Motivos decorativos sutiles
			\fill[USTBlue!10] ([xshift=1.05cm,yshift=-1.8cm]current page.north west) circle (0.72cm);
			\fill[USTTeal!12] ([xshift=1.65cm,yshift=-4.3cm]current page.north west) circle (0.42cm);
			\fill[USTBlue!6] ([xshift=-1.15cm,yshift=1.0cm]current page.south east) circle (1.15cm);
			
			% Caja central con sombra
			\fill[black!6,rounded corners=7pt]
			([xshift=-4.95cm,yshift=-3.25cm]current page.center) rectangle
			([xshift=4.95cm,yshift=2.45cm]current page.center);
			
			\fill[white,rounded corners=7pt]
			([xshift=-5.1cm,yshift=-3.1cm]current page.center) rectangle
			([xshift=4.8cm,yshift=2.6cm]current page.center);
			
			% Marca institucional superior
			\node[anchor=north east,text=white]
			at ([xshift=-0.65cm,yshift=-0.48cm]current page.north east)
			{\small\bfseries UNIVERSIDAD SANTO TOMÁS \quad TALCA};
			
			% Barra lateral decorativa
			\fill[USTTeal] ([xshift=-5.1cm,yshift=-3.1cm]current page.center) rectangle
			([xshift=-4.82cm,yshift=2.6cm]current page.center);
			
			% Contenido central
			\node[align=center] at ([xshift=0.15cm,yshift=-0.1cm]current page.center){%
				\begin{minipage}{0.66\paperwidth}
					\centering
					
					{\Large\bfseries\color{USTText}\insertauthor\par}
					
					\vspace{0.28cm}
					
					
					
					\vspace{0.95cm}
					
					{\fontsize{22}{26}\selectfont\bfseries\color{USTBlueDark}\inserttitle\par}
					
					\vspace{0.28cm}
					
					{\large\color{USTTeal}\insertsubtitle\par}
					
					\vspace{0.95cm}
					
					{\small\color{USTGray}\insertdate\par}
				\end{minipage}
			};
			
			% Línea inferior fina
			\draw[USTGray!35,line width=0.35pt]
			([xshift=0.95cm,yshift=0.72cm]current page.south west) --
			([xshift=-0.95cm,yshift=0.72cm]current page.south east);
			
			% Pie inferior
			\node[anchor=south,text=USTGray]
			at ([yshift=0.2cm]current page.south)
			{\scriptsize
				\textbf{Fundamentos de Matemática Básica}
				\;|\;
				Agronomía
				\;|\;
				Universidad Santo Tomás Talca};
			
		\end{tikzpicture}
	\end{frame}
}

% ----------------------------------------------------------
% DATOS DEL CURSO
% ----------------------------------------------------------

\title{Fundamentos de Matemática Básica}
\subtitle{Potencias y notación científica}
\author{}
\institute{Universidad Santo Tomás}
\date{Año académico 2026}

% ==========================================================
% DOCUMENTO
% ==========================================================

\begin{document}
	
	\PortadaProfesionalUST
	
	% ----------------------------------------------------------
	\begin{frame}{Contenidos de la clase}
	
	\begin{block}{Aprendizajes de la sesión}
		\begin{enumerate}
			\item Comprender el concepto de \destacar{potencia}.
			\item Interpretar la \destacar{base} y el \destacar{exponente}.
			\item Aplicar las \destacar{propiedades de las potencias}.
			\item Reconocer \destacar{patrones numéricos}.
			\item Resolver \destacar{problemas contextualizados en Agronomía}.
		\end{enumerate}
	\end{block}
	
	\end{frame}
	
	% ----------------------------------------------------------
	\begin{frame}{Concepto de potencia}
		
		\begin{block}{Definición}
			Una \destacar{potencia} representa una multiplicación repetida de un mismo número.
			
			\[
			a^n = a\cdot a\cdot a \cdots a
			\]
			
			donde:
			
			\begin{itemize}
				\item \(a\) es la \destacar{base}
				\item \(n\) es el \destacar{exponente}
			\end{itemize}
			
		\end{block}
		
		\pause
		
		\begin{exampleblock}{Ejemplos}
			
			\[
			3^4 = 3\cdot3\cdot3\cdot3 = 81
			\]
			
			\[
			2^5 = 32
			\]
			
			\[
			10^3 = 1000
			\]
			
		\end{exampleblock}
		
	\end{frame}
	
	% ----------------------------------------------------------
	\begin{frame}{¿Por qué estudiar potencias en Agronomía?}
		\begin{columns}[T,totalwidth=\textwidth]
			\begin{column}{0.52\textwidth}
				\begin{block}{Situaciones reales}
					\begin{itemize}
						\item Número de semillas por hectárea.
						\item Cantidad de bacterias en una muestra de suelo.
						\item Concentración de micronutrientes.
						\item Conversión de unidades y escalas.
					\end{itemize}
				\end{block}
			\end{column}
			
			\begin{column}{0.44\textwidth}
				\begin{exampleblock}{Ejemplo rápido}
					Una dosis puede expresarse como:
					\[
					0{,}00045\ \text{kg} = 4{,}5\times 10^{-4}\ \text{kg}
					\]
					La segunda forma es más clara para calcular, comparar y comunicar.
				\end{exampleblock}
			\end{column}
		\end{columns}
		
		\vspace{0.2cm}
		\centering
		{\small \textcolor{WarmGray}{La matemática permite simplificar, modelar y validar información técnica.}}
	\end{frame}
	
	% ----------------------------------------------------------
	\begin{frame}{Concepto de potencia}
		\begin{block}{Definición}
			Una \destacar{potencia} representa una multiplicación repetida de un mismo factor:
			\[
			a^n = \underbrace{a\cdot a\cdot a\cdots a}_{n \text{ veces}}
			\]
			donde:
			\begin{itemize}
				\item \(a\) es la \destacar{base},
				\item \(n\) es el \destacar{exponente}.
			\end{itemize}
		\end{block}
		
		\pp
		
		\begin{exampleblock}{Ejemplos}
			\[
			2^4 = 2\cdot2\cdot2\cdot2 = 16
			\qquad\qquad
			10^3 = 1000
			\]
			\[
			5^2 = 25
			\qquad
			10^{-2}=\frac{1}{10^2}=\frac{1}{100}=0{,}01
			\]
		\end{exampleblock}
	\end{frame}
	
	% ----------------------------------------------------------
	\begin{frame}{Propiedades de las potencias}
		
		\begin{block}{Reglas fundamentales}
			
			\[
			a^m \cdot a^n = a^{m+n}
			\]
			
			\[
			\frac{a^m}{a^n} = a^{m-n}
			\]
			
			\[
			(a^m)^n = a^{mn}
			\]
			
			\[
			(ab)^n = a^n b^n
			\]
			
		\end{block}
		
		\pause
		
		\begin{alertblock}{Importante}
			
			\[
			a^m + a^n \neq a^{m+n}
			\]
			
		\end{alertblock}
		
	\end{frame}
	
	% ----------------------------------------------------------
	\begin{frame}{Patrones y relaciones en las potencias de 10}
		\begin{columns}[T,totalwidth=\textwidth]
			\begin{column}{0.48\textwidth}
				\begin{block}{Observación}
					\[
					10^1 = 10
					\]
					\[
					10^2 = 100
					\]
					\[
					10^3 = 1000
					\]
					\[
					10^4 = 10000
					\]
					Cada vez que el exponente aumenta en 1, el número se multiplica por 10.
				\end{block}
			\end{column}
			
			\begin{column}{0.48\textwidth}
				\begin{exampleblock}{Hacia exponentes negativos}
					\[
					10^0 = 1
					\]
					\[
					10^{-1} = 0{,}1
					\]
					\[
					10^{-2} = 0{,}01
					\]
					\[
					10^{-3} = 0{,}001
					\]
					Cada vez que el exponente disminuye en 1, el número se divide por 10.
				\end{exampleblock}
			\end{column}
		\end{columns}
		
		\vspace{0.2cm}
		\begin{ideaBox}
			\textbf{Patrón clave:} los exponentes permiten describir regularidades de crecimiento o disminución en escala decimal.
		\end{ideaBox}
	\end{frame}
	
	% ----------------------------------------------------------
	\begin{frame}{Ejercicio contextualizado}
		
		\begin{block}{Situación agrícola}
			
			Una parcela experimental utiliza:
			
			\[
			2^5
			\]
			
			semillas por metro cuadrado.
			
			Si el cultivo cubre:
			
			\[
			2^3
			\]
			
			metros cuadrados, ¿cuántas semillas se utilizan en total?
			
		\end{block}
		
		\pause
		
		\begin{exampleblock}{Resolución}
			
			\[
			2^5 \cdot 2^3 = 2^{5+3} = 2^8
			\]
			
			\[
			2^8 = 256
			\]
			
			Se utilizan aproximadamente **256 semillas**.
			
		\end{exampleblock}
		
	\end{frame}
	
	% ----------------------------------------------------------
	\begin{frame}{Crecimiento de bacterias en el suelo}
		
		\begin{block}{Contexto}
			
			Una colonia bacteriana se duplica cada hora.
			
			La cantidad inicial es:
			
			\[
			2^3
			\]
			
			bacterias.
			
			Después de 4 horas:
			
			\[
			2^3 \cdot 2^4
			\]
			
		\end{block}
		
		\pause
		
		\begin{exampleblock}{Cálculo}
			
			\[
			2^3 \cdot 2^4 = 2^{7}
			\]
			
			\[
			2^7 = 128
			\]
			
			Después de 4 horas hay **128 bacterias**.
			
		\end{exampleblock}
		
	\end{frame}
	
	% ----------------------------------------------------------
	\begin{frame}{Fertilización}
		
		\begin{block}{Situación}
			
			Un fertilizante se aplica en dosis proporcionales a potencias de 10.
			
			En un ensayo se utiliza:
			
			\[
			5\times10^2
			\]
			
			gramos por hectárea.
			
			Si el cultivo ocupa:
			
			\[
			10^3
			\]
			
			hectáreas experimentales, calcule la cantidad total.
			
		\end{block}
		
		\pause
		
		\begin{exampleblock}{Resolución}
			
			\[
			(5\times10^2)(10^3)
			\]
			
			\[
			=5\times10^{5}
			\]
			
			Se requieren aproximadamente:
			
			\[
			500000\text{ gramos}
			\]
			
		\end{exampleblock}
		
	\end{frame}
	% ----------------------------------------------------------
	\begin{frame}{Fertilización}
		
		\begin{block}{Situación}
			
			Un fertilizante se aplica en dosis proporcionales a potencias de 10.
			
			En un ensayo se utiliza:
			
			\[
			5\times10^2
			\]
			
			gramos por hectárea.
			
			Si el cultivo ocupa:
			
			\[
			10^3
			\]
			
			hectáreas experimentales, calcule la cantidad total.
			
		\end{block}
		
		\pause
		
		\begin{exampleblock}{Resolución}
			
			\[
			(5\times10^2)(10^3)
			\]
			
			\[
			=5\times10^{5}
			\]
			
			Se requieren aproximadamente:
			
			\[
			500000\text{ gramos}
			\]
			
		\end{exampleblock}
		
	\end{frame}
	
	% ----------------------------------------------------------
	\begin{frame}{Ejercicios}
		
		Resuelva:
		
		\begin{enumerate}
			
			\item \(2^4 \cdot 2^3\)
			
			\item \((3^2)^3\)
			
			\item \(5^3 \cdot 5^2\)
			
			\item \(\dfrac{10^6}{10^2}\)
			
			\item En un cultivo se plantan \(3^4\) semillas por parcela.
			Si hay \(3^2\) parcelas, ¿cuántas semillas hay en total?
			
		\end{enumerate}
		
	\end{frame}
	
	% ----------------------------------------------------------
	\begin{frame}{Patrones}
		
		Observe la secuencia:
		
		\[
		2^1 = 2
		\]
		
		\[
		2^2 = 4
		\]
		
		\[
		2^3 = 8
		\]
		
		\[
		2^4 = 16
		\]
		
		\begin{block}{Preguntas}
			
			\begin{itemize}
				
				\item ¿Qué ocurre cuando el exponente aumenta en 1?
				
				\item ¿Cómo crecerá la secuencia?
				
			\end{itemize}
			
		\end{block}
		
	\end{frame}
	
	% ----------------------------------------------------------
\begin{frame}{Síntesis}
	
	\begin{block}{Ideas clave}
		
		\begin{itemize}
			
			\item Las potencias representan multiplicaciones repetidas.
			
			\item Las propiedades de las potencias permiten simplificar cálculos.
			
			\item Las potencias aparecen frecuentemente en fenómenos de crecimiento.
			
			\item En Agronomía permiten modelar poblaciones, concentraciones y escalas.
			
		\end{itemize}
		
	\end{block}
	
\end{frame}	
	% ----------------------------------------------------------
	\begin{frame}{Trabajo personal del estudiante}
		\begin{block}{Ejercicios propuestos}
			Resuelva en su cuaderno:
			\begin{enumerate}
				\item Exprese en notación científica:
				\[
				580000,\qquad 0{,}00073,\qquad 91200000
				\]
				\item Simplifique:
				\[
				2^3\cdot2^5,\qquad \frac{10^7}{10^3},\qquad (3^2)^4
				\]
				\item Calcule:
				\[
				(4\times10^3)(2\times10^5)
				\qquad
				\frac{9\times10^6}{3\times10^2}
				\]
				\item Explique el error en:
				\[
				0{,}034 = 34\times10^{-3}
				\]
				¿Es equivalente? ¿Está en notación científica?
			\end{enumerate}
		\end{block}
		
		\begin{exampleblock}{Proyección}
			En la siguiente clase se aplicarán estas herramientas a conversiones, escalas y análisis de datos cuantitativos del ámbito agrícola.
		\end{exampleblock}
	\end{frame}
	
\end{document}